In this appendix we provide a short script which gives an alternative way of describing an explicit expansion of the cocycle $T_5|T_{6} + d(c_4)$ modulo $(8,v_1^4)$ was used in the proof (of \Cref{lem:prod-big}).
This script runs within the computer algebra system Sage and
evaluates the cobar differential in the Hopf algebroid $(\BP_*, \BP_*\BP)$ modulo $v_1$ and a power of $2$ on terms which include only $t_1$ and $t_2$.
Ultimately, we know this script gives the correct answer in our case of interest since we have checked the outputs by hand.

\begin{itemize}
\item The ring $K$ is the base ring for our computation. We use $\Z/8$ since that is what is relevant for us, but other choices work as well.
\item The rings $A$ and $B$ represent the third and fourth layers of the cobar complex respectively.
\item The modifiers $a,b,c$ or $w,x,y,z$ indicate which tensor factor a given variable comes from.
\item The cobar complex is a cosimplicial ring and the various $di$'s are relevant face maps of this cosimplicial ring.
\item The term $E$ corresponds to the cocycle $T_j|T_{j+1}$ from \Cref{sec:alg-nov-diff}.
\item The final line (which is expanded and printed on evaluation) is the cocycle $T_j|T_{j+1} + d(c_j)$. \item We have annotated the various terms which make up $c_j$ with the length of the May differential they are involved in (see \Cref{subsec:correction}).
\end{itemize}

% evaluates the cobar differential modulo $v_1$ in the presentation of $\Mfg$ via the Hopf algebroid $(\BP_*, \BP_*\BP)$.

% Our code includes only $t_1$ and $t_2$ since these are the only terms that appear in our application.


% The second snippet is a python script which takes a cocycle in the cobar complex (such as those which appears in the output of the first script) and calculates the May filtration of each monomial. We warn the reader that the second script in particular is quite delicate (for example it will break in the presence of a minus sign in the input). These scripts are provided as-is without any guarantee that they can be used beyond their original intended examples. Ultimately, we know these scripts give correct answers in our cases of interest since we have checked the outputs by hand. 

% \subsection{cobar.sage}\ 

% This sage script computes the cobar differential in the Hopf algebroid $(\BP_*, \BP_*\BP)$ modulo $v_1$ on terms which include only $t_1$ and $t_2$. The ring $K$ is the base ring for our computation. We use $\Z/8$ since that is what is relevant for us, but other choices work as well. The rings $A$ and $B$ represent the third and fourth layers of the cobar complex respectively. The modifiers $a,b,c$ or $w,x,y,z$ indicate which tensor factor a given variable comes from. The cobar complex is a cosimplicial ring and the various $di$'s are relevant face maps of this cosimplicial ring. The term $E$ corresponds to the cocyle $T_j|T_{j+1}$ from \Cref{sec:alg-nov-diff}. The final line (which is expanded and printed on evaluation) is the cocycle $T_j|T_{j+1} + d(c_j)$. We have annotated the various terms which make up $c_j$ with the length of the May differential they are involved in (see \Cref{subsec:correction}).

\begin{verbatim}

K = Zmod(8);
A.<t1a, t1b, t1c, t2a, t2b, t2c> = K[];
B.<t1w, t1x, t1y, t1z, t2w, t2x, t2y, t2z> = K[];

d0 = A.hom([t1x, t1y, t1z, t2x, t2y, t2z], B);
d1 = A.hom([t1w + t1x, t1y, t1z, t2w - t1w * t1x^2 + t2x, t2y, t2z], B);
d2 = A.hom([t1w, t1x + t1y, t1z, t2w, t2x - t1x * t1y^2 + t2y, t2z], B);
d3 = A.hom([t1w, t1x, t1y + t1z, t2w, t2x, t2y - t1y * t1z^2 + t2z], B);
d4 = A.hom([t1w, t1x, t1y, t2w, t2x, t2y], B);

def cobar_diff(x):
    return d0(x) - d1(x) + d2(x) - d3(x) + d4(x);

def make_theta_left(k):
    return 4 * t1w^(1 * 2^(k)) * t1x^(7 * 2^(k)) 
         + 6 * t1w^(2 * 2^(k)) * t1x^(6 * 2^(k)) 
         + 4 * t1w^(3 * 2^(k)) * t1x^(5 * 2^(k)) 
         + 3 * t1w^(4 * 2^(k)) * t1x^(4 * 2^(k)) 
         + 4 * t1w^(5 * 2^(k)) * t1x^(3 * 2^(k)) 
         + 6 * t1w^(6 * 2^(k)) * t1x^(2 * 2^(k)) 
         + 4 * t1w^(7 * 2^(k)) * t1x^(1 * 2^(k));

def make_theta_right(k):
    return 4 * t1y^(1 * 2^(k)) * t1z^(7 * 2^(k)) 
         + 6 * t1y^(2 * 2^(k)) * t1z^(6 * 2^(k)) 
         + 4 * t1y^(3 * 2^(k)) * t1z^(5 * 2^(k)) 
         + 3 * t1y^(4 * 2^(k)) * t1z^(4 * 2^(k)) 
         + 4 * t1y^(5 * 2^(k)) * t1z^(3 * 2^(k)) 
         + 6 * t1y^(6 * 2^(k)) * t1z^(2 * 2^(k)) 
         + 4 * t1y^(7 * 2^(k)) * t1z^(1 * 2^(k));

# The n which appears here is related to j by n = j - 3.
n = 2;

theta_left = make_theta_left(n);
theta_right = make_theta_right(n+1);

E = theta_left * theta_right;

term1 = t1a^(4 * 2^n) * t1b^(12 * 2^n) * t1c^(8 * 2^n);
term2 = 7 * t2a^(4 * 2^n) * t1b^(4 * 2^n) * t1c^(8 * 2^n);
term3 = 2 * t2a^(2 * 2^n) * t2b^(2 * 2^n) * t2c^(4 * 2^n);
term4 = 2 * t1a^(2 * 2^n) * t1b^(4 * 2^n) * t2c^(4 * 2^n) 
          * (t2a^(2 * 2^n) + t2b^(2 * 2^n));
term5a = 2 * t1a^(4 * 2^n) * t1b^(4 * 2^n) * t1c^(4 * 2^n);
term5b = 2 * t1a^(2 * 2^n) * t1b^(2 * 2^n) * t1c^(8 * 2^n);
term5 = (term5a + term5b) * (t2a^(4 * 2^n) + t2b^(4 * 2^n) + t2c^(4 * 2^n));
term6 = 2 * t1a^(6 * 2^n) * t1b^(10 * 2^n) * t1c^(8 * 2^n);
term7 = 2 * t1a^(8 * 2^n) * t1b^(12 * 2^n) * t1c^(4 * 2^n) 
      + 2 * t1a^(12 * 2^n) * t1b^(8 * 2^n) * t1c^(4 * 2^n); 
term8 = 2 * t1a^(2 * 2^n) * t2b^(2 * 2^n) * t1c^(16 * 2^n) 
      + 2 * t1a^(4 * 2^n) * t2b^(4 * 2^n) * t1c^(8 * 2^n);

correction = term1 + term2; # first row
correction += term3; # may d1
correction += term4; # may d0
correction += term5; # may d0
correction += term6; # may d0
correction += term7; # may d0
correction += term8; # may d1
E + cobar_diff(correction);

\end{verbatim}

% \subsection{may\_fil.py}\

% This short python script can be used to take a string which represents an element of the cobar complex and to split it into monomials and evaluate the May filtration of each monomoial. Our method of splitting things up is not particularly robust and this script only really works well when each monomial appears with coefficient ``2''.

% \begin{verbatim}
% def dyadic_digit_sum(n):
%     if n == 0:
%         return 0
%     return (n%2) + dyadic_digit_sum(n/2)

% def term_to_vec(term):
%     term_list = term.split("*")
%     vec = dict()
%     for t in term_list[1:]:
%         key, val = t.split("^")
%         key = key.strip()
%         val = (2 * int(key[1]) - 1) * dyadic_digit_sum(int(val))
%         vec[key] = val
%     return vec

% def vec_may_fil(vec):
%     m = 0
%     for val in vec.values():
%         m += val
%     return m

% def may_fil(string):
%     term_list = string.split("+")
%     return [ (t, vec_may_fil(term_to_vec(t))) for t in term_list ]

% def top_may_fil(may_fil_list):
%     fils = [ m for (t,m) in may_fil_list ]
%     M = max(fils)
%     return (M, [ t for (t,m) in may_fil_list if m == M] )
% \end{verbatim}


% \begin{verbatim}

% K = Zmod(8);
% A.<t1a, t1b, t1c, t2a, t2b, t2c> = K[];
% B.<t1w, t1x, t1y, t1z, t2w, t2x, t2y, t2z> = K[];

% d0 = A.hom([t1x, t1y, t1z, t2x, t2y, t2z], B);
% d1 = A.hom([t1w + t1x, t1y, t1z, t2w + t1w^2 * t1x + t2x, t2y, t2z], B);
% d2 = A.hom([t1w, t1x + t1y, t1z, t2w, t2x + t1x^2 * t1y + t2y, t2z], B);
% d3 = A.hom([t1w, t1x, t1y + t1z, t2w, t2x, t2y + t1y^2 * t1z + t2z], B);
% d4 = A.hom([t1w, t1x, t1y, t2w, t2x, t2y], B);

% def cobar_diff(x):
%     return d0(x) - d1(x) + d2(x) - d3(x) + d4(x)

% n = 2
% theta_left =   4 * t1w^(1 * 2^(n+1)) * t1x^(7 * 2^(n+1))
%              + 6 * t1w^(2 * 2^(n+1)) * t1x^(6 * 2^(n+1))
%              + 4 * t1w^(3 * 2^(n+1)) * t1x^(5 * 2^(n+1))
%              + 3 * t1w^(4 * 2^(n+1)) * t1x^(4 * 2^(n+1))
%              + 4 * t1w^(5 * 2^(n+1)) * t1x^(3 * 2^(n+1))
%              + 6 * t1w^(6 * 2^(n+1)) * t1x^(2 * 2^(n+1))
%              + 4 * t1w^(7 * 2^(n+1)) * t1x^(1 * 2^(n+1))
% theta_right =   4 * t1y^(1 * 2^n) * t1z^(7 * 2^n)
%              + 6 * t1y^(2 * 2^n) * t1z^(6 * 2^n)
%              + 4 * t1y^(3 * 2^n) * t1z^(5 * 2^n)
%              + 3 * t1y^(4 * 2^n) * t1z^(4 * 2^n)
%              + 4 * t1y^(5 * 2^n) * t1z^(3 * 2^n)
%              + 6 * t1y^(6 * 2^n) * t1z^(2 * 2^n)
%              + 4 * t1y^(7 * 2^n) * t1z^(1 * 2^n)


% E = -1 * theta_left * theta_right

% term1 = t1a^(8 * 2^n) * t1b^(12 * 2^n) * t1c^(4 * 2^n)
% term2 = 7 * t1a^(8 * 2^n) * t1b^(4 * 2^n) * t2c^(4 * 2^n)
% term3 = 2 * t1a^(4 * 2^n) * t1b^(4 * 2^n) * t1c^(4 * 2^n) * t2a^(4 * 2^n)
% term4 = 2 * t1a^(4 * 2^n) * t1b^(4 * 2^n) * t1c^(4 * 2^n) * t2b^(4 * 2^n)
% term5 = 2 * t1a^(4 * 2^n) * t1b^(4 * 2^n) * t1c^(4 * 2^n) * t2c^(4 * 2^n)
% term6 = 2 * t1a^(4 * 2^n) * t1b^(8 * 2^n) * t2c^(4 * 2^n)
% term7 = 2 * t1a^(4 * 2^n) * t2b^(4 * 2^n) * t1c^(8 * 2^n)
% term8 = 2 * t1a^(4 * 2^n) * t1b^(12 * 2^n) * t1c^(8 * 2^n)
% term9  = 2 * t2a^(4 * 2^n) * t1b^(4 * 2^n) * t2b^(2 * 2^n) * t1c^(2 * 2^n)
% term10 = 2 * t2a^(4 * 2^n) * t1b^(4 * 2^n) * t2c^(2 * 2^n) * t1c^(2 * 2^n)
% term11 = 2 * t2a^(4 * 2^n) * t2b^(2 * 2^n) * t2c^(2 * 2^n)
% term12 = 2 * t1a^(8 * 2^n) * t2a^(4 * 2^n) * t1b^(2 * 2^n) * t1c^(2 * 2^n)
% term13 = 2 * t1a^(16 * 2^n) * t2b^(2 * 2^n) * t1c^(2 * 2^n)
% term14 = 2 * t1a^(8 * 2^n) * t1b^(2 * 2^n) * t2b^(4 * 2^n) * t1c^(2 * 2^n)
% term15 = 2 * t1a^(8 * 2^n) * t1b^(2 * 2^n) * t1c^(2 * 2^n) * t2c^(4 * 2^n)
% term16 = 2 * t1a^(8 * 2^n) * t1b^(10 * 2^n) * t1c^(6 * 2^n)

% correction = term1 + term2 + term3 + term4 + term5 + term6 
%            + term7 + term8 + term9 + term10 + term11 + term12 
%            + term13 + term14 + term15 + term16

% E + cobar_diff(correction)

% \end{verbatim}


%%% Local Variables:
%%% mode: latex
%%% TeX-master: "main"
%%% End:
